\chapter{Appendix Theta: The Foundry Scenario}

\section{The Legacy Dependency (Technical \\\& Cultural)}
The Theta Foundry relies heavily on the state power grid (industrial energy consumption), legacy supply chains for raw materials (steel, silicon, polymers), and state-mandated industrial safety certifications (OSHA/ISO). Culturally, it depends on a specialized, specialized labor force accustomed to traditional corporate hierarchies and intellectual property monopolies.

\section{The Parasitic Bridge (Technical \\\& Legal)}
Layer 7 Physical Grid Leeching is aggressively employed. The foundry draws from the legacy macro-grid during off-peak hours when fiat energy costs are lowest, storing energy in on-site Layer 2 battery arrays. Raw materials are sourced via Web2 API integration with legacy logistics firms. The operation is shielded by a corporate LLC that holds the necessary ISO certifications, while internally operating as an open-source, decentralized manufacturing collective.

\section{The Decoupling Trigger}
Decoupling from the state macro-grid occurs when on-site renewable generation (solar/geothermal) and battery buffering reach 150\% of the foundry's peak load requirement. Decoupling from legacy supply chains occurs as the Sovereign Stack's internal circular economy (recycling and localized resource extraction) provides 80\% of required raw inputs.

\section{Orchestration \\\& Ethical Mechanics}
The orchestration uses continuous API monitoring of energy markets to switch between macro-grid leeching and islanded operation. Ethically, the foundry must not become a source of pollution for marginalized communities; the Layer 7 bridge includes strict, transparent cryptographic auditing of emissions and waste, exceeding state regulations and enforcing accountability via DAO slashing mechanisms.
